\section{Centering at the layer cuts}\label{sec:centering}

We subtract independent kinetic factors while keeping each nonsingleton
interaction component intact. Correlation errors obtained by removing
independent factors were developed in~\cite{PulvirentiSimonella2017};
their propagation across time layers with the earlier histories retained
is central to~\cite{DHM25}. This section proves the componentwise
centering identities needed for the present potential class. We obtain
an exact recurrence across the layers, then show that each surviving
centered root must lead to a small input or to a contact that prevents
its family from evolving independently.

\subsection{The patternwise transform}
For a partition $\pi$ of a finite label set $R$, write
$\chi_\pi^\eps(Z_R)$ for the indicator that the connected components of
$\{\{i,j\}:|x_i-x_j|\leq\eps\}$ are exactly the blocks of $\pi$.
Let $S(\pi)$ be the set of labels in singleton blocks and let
$B(\pi)=R\setminus S(\pi)$. For $A\subset S(\pi)$, the partition
$\pi_A$ consists of the nonsingleton blocks of $\pi$ and the singletons
in $A$; its label set is $B(\pi)\cup A$. We use
$f_D(t,Z_D)=\prod_{i\in D}f(t,z_i)$, with $f_\varnothing=1$.

\begin{definition}[Patternwise centered marginal]\label{cent:definition}
The centered marginal on pattern $\pi$ is
\begin{equation}\label{cent:definition-eq}
 E_\pi^\eps(t,Z_R)=\chi_\pi^\eps(Z_R)
 \sum_{A\subset S(\pi)}(-1)^{|S(\pi)\setminus A|}
 f_{S(\pi)\setminus A}(t,Z_{S(\pi)\setminus A})
 \widetilde F_{B(\pi)\cup A}^\eps(t,Z_{B(\pi)\cup A}).
\end{equation}
We set $E_\varnothing^\eps=1$. A nonsingleton block is always retained
as a whole; its individual labels are not centered.
\end{definition}

The indicator in \eqref{cent:definition-eq} belongs to the full array.
For example, deleting a singleton $i$ from a lower marginal does not
remove the conditions $|x_i-x_j|>\eps$ for the other displayed labels.
These conditions multiply the lower marginal and the independent
factor $f(t,z_i)$.

\begin{proposition}[Centering and inversion]\label{cent:identity}
For every $\pi$ and almost every $Z_R$,
\begin{equation}\label{cent:inverse}
 \chi_\pi^\eps\widetilde F_R^\eps(t)
 =\chi_\pi^\eps
 \sum_{A\subset S(\pi)}f_{S(\pi)\setminus A}(t)
 E_{\pi_A}^\eps(t).
\end{equation}
Both formulas hold for arbitrary factorial densities, without a
smallness or factorization assumption.
\end{proposition}
\begin{proof}
On $\{\chi_\pi^\eps=1\}$ every restricted pattern $\pi_A$ has indicator
one. Substitute \eqref{cent:definition-eq} on the right of
\eqref{cent:inverse}. For $D\subset S(\pi)$, the coefficient of
$f_{S(\pi)\setminus D}\widetilde F_{B(\pi)\cup D}^\eps$ is
\[
 \sum_{D\subset A\subset S(\pi)}(-1)^{|A\setminus D|}
 =(1-1)^{|S(\pi)\setminus D|}.
\]
Only $D=S(\pi)$ survives. Off that pattern both sides vanish.
\end{proof}

\begin{corollary}[Collapsing the pattern sum]\label{cent:pattern-collapse}
For $D\subset R$ let
\begin{equation}\label{cent:singleton-deletion}
 \iota_D^R(Z_R)
 =\prod_{i\in D}\prod_{j\in R\setminus\{i\}}
       \1_{\{|x_i-x_j|>\eps\}},\qquad
 \mathcal E_R=\sum_{\pi\in\mathcal P(R)}E_\pi,
\end{equation}
where $\mathcal P(R)$ is the set of partitions of $R$. Then
\begin{align}
 \mathcal E_R
 &=\sum_{D\subset R}(-1)^{|D|}\iota_D^R f_D
                       \widetilde F_{R\setminus D},
       \label{cent:collapsed-centering}\\
 \widetilde F_R
 &=\sum_{D\subset R}\iota_D^R f_D
                       \mathcal E_{R\setminus D}.
       \label{cent:collapsed-inverse}
\end{align}
Thus a cut carries subset choices, but no additional factor equal to
the number of partitions. For arbitrary signed families $G$ and $e$
with each $e_\pi$ supported on its pattern, write
\[
 \|G\|_{1,\leq U}=\max_{0\leq r\leq U}\|G_{[r]}\|_1,
 \qquad
 \|e\|_{\pi,1,\leq U}
   =\max_{0\leq r\leq U}\sum_{\pi\in\mathcal P([r])}\|e_\pi\|_1.
\]
Let $\mathsf M_t$ be the transform \eqref{cent:definition-eq} and
$\mathsf I_t$ its inverse \eqref{cent:inverse}, summed over full
patterns. They satisfy
\begin{equation}\label{cent:centering-norm}
 \|\mathsf M_tG\|_{\pi,1,\leq U}\leq2^U\|G\|_{1,\leq U},
 \qquad
 \|\mathsf I_te\|_{1,\leq U}\leq2^U\|e\|_{\pi,1,\leq U}.
\end{equation}
\end{corollary}
\begin{proof}
At a fixed array exactly one full pattern occurs. Its singleton
labels are exactly the labels which may belong to a set $D$ with
$\iota_D^R=1$. Deleting these labels leaves the unique restricted
pattern. Summing \eqref{cent:definition-eq} and
\eqref{cent:inverse} therefore gives
\eqref{cent:collapsed-centering}--\eqref{cent:collapsed-inverse}.
For disjoint $D,A$ one can also check inversion directly from
\[
 \iota_D^R\,\iota_A^{R\setminus D}=\iota_{D\cup A}^R
\]
and the binomial cancellation in Proposition~\ref{cent:identity}.
Because pattern supports are disjoint, the sum of the patternwise
absolute values is the absolute value of their assembled density.
Taking absolute values in the collapsed formulas, dropping
$\iota_D^R\leq1$, and integrating each deleted factor using
$\int f(t)=1$ proves \eqref{cent:centering-norm}.
\end{proof}

At time zero, the insertion formula of Proposition~\ref{pre:initial}
gives $E_\pi^\eps(0)=0$ if $B(\pi)\ne\varnothing$: every term in
\eqref{cent:definition-eq} retains all nonsingleton blocks, hence
contains a close pair on the full pattern and has zero initial
factorial density. The full-array indicator remains in place when
singleton labels are deleted. For a singleton pattern with $k$ labels,
the separated-array estimate of that proposition gives
\begin{equation}\label{cent:initial-use}
 |E_\pi^\eps(0,Z_R)|\leq(C\eps^c)^k f_{0,R}(Z_R).
\end{equation}
In particular, this small initial factor is attached to each retained
singleton slot, rather than to each layer through which it is carried.

\subsection{An explicit one-root reference sum}

We next construct the signed binary histories against which an
independently evolving singleton family will cancel. Their construction
determines the exact pairing in Proposition~\ref{cent:cancellation}.
The separate comparison with $f(t)$ bounds the residual left by this
replacement at one fixed Gaussian exponent.

Fix $u<t$ with $t-u\leq b$. Use the cells of this layer, of length at
most $b/J$. In the following construction $z_1(t)=z$ is fixed. A tree
with $n$ births has chronological backward labels $2,\ldots,n+1$,
ordered times
\[
 t>s_1>\cdots>s_n>u,
 \qquad a_i\in\{1,\ldots,i\},\quad \sigma_i\in\{+,-\}.
\]
At $s_i$ label $i+1$ is placed at
$x_{a_i}(s_i)+\eps\omega_i$, with velocity $w_i$; write
$g_i=w_i-v_{a_i}(s_i+)$. Here $\omega_i\in\mathbb S^2$ and
$g_i\cdot\omega_i>0$. The plus choice follows the isolated two-body
flow backward from this outgoing sphere crossing until the incoming
sphere crossing, then continues the two lines freely. The minus choice
continues both lines freely through the ball. Its sign is $-1$.
All other lines move freely while this pair is being resolved.
The contact time is an integration parameter, not a force-system cut:
the new label also has its outgoing free segment from $s_i$ to the
\emph{upper boundary of its insertion cell}. Its formal lifetime ends
at that cell boundary, where it is separated from the other labels.
That outgoing segment is a displayed trajectory in every avoidance
test, for both signs. The free passage has duration
$2\eps(g_i\cdot\omega_i)/|g_i|^2$; the true passage has the duration
controlled by Lemma~\ref{pre:duration}.

Keep only parameters for which each scheduled pair passage finishes
before the earlier boundary of its insertion cell, no two births of
this tree occupy one cell, and there are no other contacts between any
of its displayed lines. At every cut at which a line starts or stops,
require separation from the other displayed lines. In particular the
array at $u$ consists of singleton blocks. Denote this indicator by
$\mathscr S_{n}^{\eps,J}$. It is zero if any scheduled passage fails to
exit. The resulting bottom array is $Z_{n+1}^{\eps}(u)$.
These are \emph{reference histories}: the free negative passage is a
term of a signed identity, not a modification of the particle process.
A one-root displayed family satisfying these conditions is called
\emph{simple}.

The measure in this representation is
\begin{equation}\label{cent:tree-measure}
 d\lambda_n
 =\prod_{i=1}^n ds_i\,dw_i\,d\omega_i\,
             (g_i\cdot\omega_i)_+ .
\end{equation}
Surface measure is the ordinary area measure on $\mathbb S^2$.
To see the Jacobian directly, let $c_i$ be the upper boundary of the
insertion cell. On its outgoing segment the new label's relative
position at $c_i$ is
\[
 x_{i+1}(c_i)-x_{a_i}(c_i)
   =\eps\omega_i+(c_i-s_i)(w_i-v_{a_i}(s_i+)).
\]
The normal component of the derivative with respect to $s_i$ is
$-g_i\cdot\omega_i$, and the two tangential derivatives have area
factor $\eps^2$. Thus the spatial Jacobian is
$\eps^2(g_i\cdot\omega_i)_+$. Restricting to the first backward
crossing makes this parametrization one-to-one, up to the null
set of tangencies. Hamiltonian volume preservation supplies the
remaining phase coordinates for the plus sign; the minus sign uses
free transport. The factor $\eps^2$ is canceled by
$\mu=\eps^{-2}$. Birth labels are in chronological order, so there is
no additional $1/n!$ in \eqref{cent:tree-measure}; the ordered time
domain supplies the factorial in estimates. Define
\begin{equation}\label{cent:simple-sum}
 \mathscr G_{u,t}^{\eps,J,K}(z)
 =\sum_{0\leq n<K-1}\ \sum_{a_1,\ldots,a_n}\
   \sum_{\sigma_1,\ldots,\sigma_n}
 \left(\prod_{i=1}^n\sigma_i\right)
 \int\mathscr S_n^{\eps,J}
       f_{n+1}(u,Z_{n+1}^{\eps}(u))\,d\lambda_n .
\end{equation}
Here $\sigma_i$ in the product means $1$ or $-1$; the $n=0$ term is
$f(u,x-(t-u)v,v)$. The number of particles is strictly less than $K$.

\begin{proposition}[Short-layer residual]\label{cent:residual}
Set $\beta_0=\beta/2$. There is $b_*>0$, depending only on the fixed
bounds in Proposition~\ref{pre:envelopes}, for which the following
holds. Suppose $b\leq b_*$, $K(\eps)\to\infty$, and
$J(\eps)\to\infty$ is a fixed power of $1+|\log\eps|$.
Then
\begin{equation}\label{cent:residual-def}
 r_{u,t}^{\eps}=f(t)-\mathscr G_{u,t}^{\eps,J,K},\qquad
 \sup_{\substack{0\leq u\leq t\leq T\\t-u\leq b}}
 \|r_{u,t}^{\eps}\|_{\mathrm{Bol},\beta_0}
 =:\eta_\eps\longrightarrow0.
\end{equation}
All the reference sums and residuals have uniformly bounded norms at
this one fixed exponent. The proof uses only short-layer tree
summability and two-body scattering.
\end{proposition}
\begin{proof}
We give the majorant and the limiting argument separately. Put
$p=n+1$ and
$\mathcal E=|v|^2+\sum_{i=1}^n|w_i|^2$. At every insertion sphere the
pair potential is zero. Elasticity consequently gives
\begin{equation}\label{cent:tree-energy}
 \sum_{k=1}^p|v_k(u)|^2=\mathcal E.
\end{equation}
Intermediate pair speeds are bounded by $C(\mathcal E+p)^{1/2}$ by
stability. These bounds hold for both signs. At a birth, summing the
flux over possible parents, with unused parent choices assigned zero,
gives
\begin{equation}\label{cent:parent-sum}
 \sum_{a\ \mathrm{available}}\int_{\mathbb S^2}
       ((w-v_a)\cdot\omega)_+\,d\omega
 \leq Cp\bigl(1+|w|^2+\mathcal E/p\bigr)^{1/2}.
\end{equation}
The inequality follows from
$\sum_a|v_a|\leq\sqrt{p\sum_a|v_a|^2}$.
It may be used successively, since its right side no longer depends
on the chosen parents or the angular variables. The product satisfies
\begin{equation}\label{cent:velocity-product}
 \prod_{i=1}^n\bigl(1+|w_i|^2+\mathcal E/p\bigr)^{1/2}
 \leq C^p(1+\mathcal E/p)^{p/2}.
\end{equation}
Indeed the arithmetic--geometric mean inequality bounds the product
by the $n$th power of its arithmetic mean, and
$\sum_i|w_i|^2\leq\mathcal E$.

Let $M$ bound both the uniform Gaussian norm and the summed spatial
Gaussian norm of $f(u)$ at exponent $\beta$. Use the spatial envelope
on the continued label $1$ and the uniform bound on all other bottom
labels. Every displacement of that label has length at most
$Cb(\mathcal E+p)^{1/2}$. Translation of its summability boxes therefore
costs at most $C(1+b\sqrt{\mathcal E+p})^3$: any unit box enlarged by
that radius is covered by this number of unit boxes. To handle the velocity supremum without interchanging it with the
spatial sum, first split the root velocities into shells
$m\leq|v|<m+1$. On each shell the displacement radius is at most
$Cb\sqrt{(m+1)^2+|W|^2+p}$, independently of the root velocity.
Sum the spatial boxes with this enlarged radius and then sum over
$m$, reserving a Gaussian factor $e^{-c m^2}$. This proves the same
bound with a convergent Gaussian shell sum in place of a root
supremum. Reserve the fixed Gaussian margin $\beta-\beta_0$. For a fixed
$\delta>0$ smaller than that margin,
\begin{equation}\label{cent:moment-bound}
 \sup_v e^{-\delta|v|^2}
 \int_{\R^{3n}} e^{-\delta\sum_i|w_i|^2}
 (1+\mathcal E/p)^{p/2}(1+b\sqrt{\mathcal E+p})^3\,dW
 \leq C^p.
\end{equation}
To verify this, use
$1+y\leq C_\delta e^{\delta y}$ with $y=\mathcal E/p$,
choose its exponent smaller than the remaining Gaussian exponent,
and absorb the fixed cubic polynomial with another fixed part of
that exponent. The resulting Gaussian integral is $C^n$; its root
supremum is finite. Its Gaussian velocity-shell version, needed in
the preceding spatial estimate, has the same $C^p$ bound. Adjusting
$C$ includes the finitely many cases $p\leq3$.

The ordered simplex has volume $(t-u)^n/n!$, and
$p^n/n!\leq e^{2p}$. Equations
\eqref{cent:tree-energy}--\eqref{cent:moment-bound}, including the
$2^n$ sign choices, prove the absolute estimate
\begin{equation}\label{cent:geometric-tail}
 \left\|\sum_{a,\sigma}\int
       \mathscr S_n^{\eps,J}
       f_p(u,Z_p^\eps(u))\,d\lambda_n\right\|_{\mathrm{Bol},\beta_0}
 \leq C_0(C_1b)^n,
\end{equation}
where the sum inside the norm is understood with absolute integrands.
Choose $b_*$ so that $C_1b_*<1/2$. The identical estimate holds for
instantaneous scattering trees.

Fix now $n$ and truncate the velocity and relative insertion parameters
to a compact set. The Gaussian majorant just proved makes the omitted
part arbitrarily small, uniformly in $u,t$ and $\eps$. Remove incoming
parameters whose unscaled residence time exceeds $D$, whose relative
speed is too small, or which belong to the null exceptional sets of
Proposition~\ref{pre:scattering}. Lemma~\ref{pre:duration} controls the
residence-time cutoff. The omitted flux tends to zero as
$D\to\infty$ and the other cutoffs are removed. On the retained set a
true or free passage has length $O(\eps D)$ and displacement
$O(\eps D)$. The set of insertion times lying within $C\eps D$ of a
cell boundary has measure at most $C_nJ\eps D$; two insertion times
in one cell cost $C_n/J$. Both quantities vanish, taking $D$ fixed
first. Shortened endpoint cells obey the same estimates.

Here is the needed null-set assertion for other meetings. Parameterize
each later birth by its velocity relative to its current parent in
laboratory axes. This is a triangular translation of the $w_i$
variables with Jacobian one. By Galilean invariance, later relative
increments do not depend on an earlier common velocity boost.
At the first divergence of two distinct rays, keep the relative
speed and the scalar scattering parameters fixed, and rotate the
relative incident velocity and impact configuration together.
The separating relative velocity has nonzero length and a spherical
direction with absolutely continuous area measure. On any subsequent
interval between scheduled births, the relative position of the two
rays has the form
\begin{equation}\label{cent:meeting-null}
 (s-s_*)\,\xi+\Gamma(s),
\end{equation}
where $s_*$ is their divergence time, $\xi$ is that relative velocity,
and $\Gamma$ is a fixed affine function after the later relative
increments and times have been fixed. There are finitely many such
intervals. For $s\ne s_*$, equality to zero confines the direction of
$\xi$ to a finite union of curves on a sphere, which has area zero.
The interval immediately after divergence is free of other scheduled
births and its nonzero relative speed separates the two rays.
For a true scattering both outgoing rays must be included when
identifying this divergence; their directions are not asserted to be
independent. This argument also covers the free sign. Fubini's theorem
therefore excludes every unprescribed meeting for fixed $n$. Relative
positions and this exceptional set are independent of the root
position and root velocity, by translation and Galilean invariance.
Thus the compact exceptional estimates are uniform for the root
supremum in the spatial norm.

Away from these exceptions, the microscopic paths converge to the
instantaneous paths, with the same velocities outside the short pair
arcs. All input positions differ by $O_n(\eps D)$, while their input
velocities agree exactly. The outgoing segments between a contact
and its upper cell boundary have length $O(b/J)$ and are also checked
for unwanted contacts. Their length tends to zero, so the same
null-set argument covers their avoidance indicators. Proposition~\ref{pre:envelopes} and the
spatial mean-value inequality give convergence of the products of
input densities in the summed spatial norm, uniformly in $u,t$.
Dominated convergence in the compact parameters, followed by removal
of the compact cutoffs, proves convergence of every fixed term.
For uniformity in the layer length, use backward elapsed times in the
fixed simplex $0<\tau_1<\cdots<\tau_n<b$ and restrict by
$\tau_n<t-u$. Extend the bottom rays freely to elapsed time $b$ in
the null-set test. This tests a larger finite time interval and still
has zero exceptional measure. The estimates for short arcs and
endpoint cells are uniform in $t-u$. If $t-u\to0$, the terms with
$n\geq1$ are directly bounded by $C_n(t-u)^n$. These observations
justify the asserted supremum over all shortened layers.

Finally the limiting insertion integrates over the outgoing
hemisphere with weight $(g\cdot\omega)_+$. For the plus sign its bottom
velocities are the inverse scattering velocities; for the minus sign
they are the unaltered velocities. The flux-preserving scattering
change of variables in Lemma~\ref{pre:operator} and the sphere-to-disk
identity \eqref{eq:pre-fluxdisk} identify their difference with precisely
$Q_\Phi$, with ordinary disk area. Iterating the mild kinetic equation
on $[u,t]$ gives the instantaneous tree sum. Its remainder tends to
zero by \eqref{cent:geometric-tail}, using the same Gaussian bound for
the terminal $f$; this also proves uniqueness on this short interval
in the class at issue. First make the common geometric tail small,
then use fixed-term convergence. This proves
\eqref{cent:residual-def} without any statement about a microscopic
large-history remainder.
\end{proof}

\subsection{Finite layer identities and cancellation}

Let $\mathsf M_t$ denote \eqref{cent:definition-eq}, acting on a family
of marginals of all orders and all patterns, and let $\mathsf I_t$
denote \eqref{cent:inverse}, summed over the full-array patterns.
They are inverse triangular transforms. Every occurrence of either
transform includes its full-array pattern indicator.
Fix a cap $U_j$ on the number of upper root labels and an integer
$K_j\geq2$. For a layer $[t_{j-1},t_j]$, compose the cell kernels
$\mathcal T^I_{R,S}$ of Proposition~\ref{cell:expansion} in backward
order and stop by the \emph{physical} total size $K_j$, as in
Proposition~\ref{cell:stopping}. Denote its completed operator by
$\mathsf C_j^{\mathrm{phys},<K_j}$ and its exact physical remainder
by $\mathsf R_j^{\mathrm{phys}}$. Before any reference-tree
replacement, the exact identity is
\begin{equation}\label{cent:physical-layer}
 \widetilde F(t_j)
 =\mathsf C_j^{\mathrm{phys},<K_j}\widetilde F(t_{j-1})
   +\mathsf R_j^{\mathrm{phys}}.
\end{equation}
Only this physical remainder is associated with the stopped-cloud
functional of Proposition~\ref{cell:stopping}.
For clarity, a completed physical term with successive arrays
$R_0,\ldots,R_J$ has pairing against an endpoint test $\psi$ equal to
\begin{equation}\label{cent:layer-kernel}
 \prod_{c=1}^J\frac{\mu^{|R_c\setminus R_{c-1}|}}
                         {|R_c\setminus R_{c-1}|!}
 \int \psi(Z_{R_0})
 \prod_{c=1}^J
 \mathsf K^{I_c}_{R_{c-1},R_c}
          (Z_{R_c};dZ_{R_{c-1}})\,G_{R_J}(Z_{R_J})\,dZ_{R_J}.
\end{equation}
The cell order is backward from the observation cut; intermediate
integrations are kernel composition. Every factor $\mathsf K$ is
its finite signed generation sum with all indicators retained.

Apply the exact transforms to \eqref{cent:physical-layer}. In the
resulting finite expression
\[
 \widehat{\mathsf A}_j
   =\mathsf M_{t_j}\mathsf C_j^{\mathrm{phys},<K_j}
                               \mathsf I_{t_{j-1}},
\]
replace each upper $f(t_j)$ factor by
$\mathscr G_{t_{j-1},t_j}^{\eps,J,K_j}+r_{t_{j-1},t_j}^{\eps}$.
Use fresh labels for the nonroot labels of each reference tree; its
root \emph{reuses the deleted upper label}. A residual is a one-body
factor on that upper label and has no trajectory inside the layer.
The lower inverse transform merely assigns existing lower labels to
independent $f(t_{j-1})$ factors or to retained centered slots.

For each term of this expansion let $p$ count all distinct formal
label lifetimes in the upper block: physical labels, reference-tree
labels, and residual upper labels, counting a reused root once.
Let $\mathsf P_{<K_j}$ and $\mathsf P_{\geq K_j}$ denote the
projections that retain the terms with the indicated value of $p$.
These are projections on the finite expanded list, not on the
microscopic configuration. Define
\begin{equation}\label{cent:reference-mismatch}
 \mathsf A_j:=\mathsf P_{<K_j}\widehat{\mathsf A}_j,
 \qquad
 B_j:=\mathsf P_{\geq K_j}\widehat{\mathsf A}_j,
 \qquad \widehat{\mathsf A}_j=\mathsf A_j+B_j.
\end{equation}
Here $B_j$ is a finite signed linear operator on the lower centered
family, with fixed independent reference and residual factors.
It is the \emph{reference cutoff mismatch}, and is distinct from
$\mathsf R_j^{\mathrm{phys}}$.

To record a uniform size bound, let $r\leq U_j$ be the full upper
order and let $d\leq r$ upper singleton slots be deleted by
$\mathsf M_{t_j}$. The physical block has $q<K_j$ labels. Each of
the $d$ deleted slots contributes either a reference tree with fewer
than $K_j$ labels, including its reused root, or one residual label.
Consequently every term in $B_j$ has
\begin{equation}\label{cent:mismatch-size}
 K_j\leq p\leq (K_j-1)+d(K_j-1)
             \leq(U_j+1)(K_j-1)<(U_j+1)K_j,
 \qquad h_0=r\leq U_j.
\end{equation}
The case of no physical root uses $q=0$ and satisfies the same bound.
No reference root creates an additional restarting top.

The exact centered recurrence is now
\begin{equation}\label{cent:layer-identity}
 E(t_j)=\mathsf A_jE(t_{j-1})
       +\mathsf M_{t_j}\mathsf R_j^{\mathrm{phys}}
       +B_j[E(t_{j-1})].
\end{equation}
Its finite iteration gives
\begin{align}\label{cent:finite-layers}
 E(t_L)={}&\mathsf A_L\cdots\mathsf A_1E(0)\nonumber\\
 &+\sum_{j=1}^L\mathsf A_L\cdots\mathsf A_{j+1}
 \left\{\mathsf M_{t_j}\mathsf R_j^{\mathrm{phys}}
                     +B_j[E(t_{j-1})]\right\}.
\end{align}
Both remainder types remain in this identity. In particular the
input of $B_j$ is the true lower centered family. Its estimate in
Section~\ref{sec:closure} expands that input only in layers strictly
earlier than $j$, treating this finite mixed upper block as
uncentered. It does not use positive stopped-cloud domination for
$B_j$, or assume an estimate of its own layer's remainder.
All these identities can first be read with the finite
particle-number restriction of Corollary~\ref{cor:cell-factorial}.

\begin{proposition}[Exact simple-family cancellation]\label{cent:cancellation}
In the retained operator $\mathsf A_j$ of
\eqref{cent:reference-mismatch}, a singleton root
whose displayed family satisfies the reference-history conditions
$\mathscr S_n^{\eps,J}$ above, ends entirely on independent
$f(t_{j-1})$, and avoids every displayed trajectory of every other
family throughout the layer cancels with its tensor-subtraction
term. The avoidance requirement includes all detached generation
groups and all reference trees coming from tensor subtraction.
It also includes the full-array pattern restrictions at cuts.
A residual factor is treated as a stopped slot, not as an unexamined
trajectory. The pairing is exact coefficientwise for total formal
size $p<K_j$. The finite complementary operator $B_j$ stays in
\eqref{cent:layer-identity}; no cancellation or positive physical
remainder bound is asserted for it.
\end{proposition}
\begin{proof}
Retain a tag saying whether each centered singleton root belongs to
the microscopic marginal in $\mathsf M_{t_j}$ or to its negative
$f(t_j)$ factor. In either case expand a simple candidate family all
the way to the lower cut and label every line it displays. Split the
domain of every diagram into the event of no cross contact with this
family and its complement.

On the no-cross-contact event, Proposition~\ref{cell:factorization}
splits the cell generation diagram into its restriction to that
family and its restriction to the other families. Explicitly, the
root subsystem flow splits; a detached connected group cannot meet
two families; each generation is the disjoint union of its family
generations. Deleting empty terminal generations in each restriction
and padding them when merging are inverse operations. The signs are
the products of the signs of the two restrictions, and the
within-generation avoidance tests split. This statement uses the
trajectories in the displayed signed diagram, including the detached
ones. Avoidance merely in the full interacting system would not
justify any of these operations.

The one-label rule in each occupied cell is exactly the plus/minus
insertion used in \eqref{cent:tree-measure}. Empty cells propagate
freely. Thus the family restriction is one of the terms of
\eqref{cent:simple-sum}, with identical times, sphere normals,
velocities, indicators, and input factors on the two tagged sides.
The inverse transform at the lower cut supplies exactly the
independent $f$ factors on its leaves. Restoring the full pattern
indicators makes the deletion and reinsertion operations inverse
also at the cuts. In particular, a nonsingleton polymer cannot enter
this pairing.

If $m$ extra labels split into $r$ in this family and $m-r$ in the
rest, the label choice contributes $\binom mr$ and
\[
 \frac{\mu^m}{m!}\binom mr
 =\frac{\mu^r}{r!}\frac{\mu^{m-r}}{(m-r)!}.
\]
Applying this in the occupied cells proves equality of the measures
and all factorial coefficients. Change only the root tag. The
centering sign reverses and every other factor stays the same.
The first eligible root in a fixed order defines an involution:
eligibility depends on the displayed paths and lower retained slots,
which the tag change preserves. The paired contributions sum to zero.

Perform the pairing coefficient by coefficient for $p<K_j$, counting
all formal lifetimes and reusing the upper root label when its tag
changes. Both the resulting physical block and each individual
reference tree then have size less than $K_j$. Hence both tagged
terms occur in the separately truncated expansion, and both belong
to $\mathsf A_j$. The pairing preserves $p$ and is an exact
involution on this retained part. The complementary terms are
precisely the already defined finite operator $B_j$, with the size
bound \eqref{cent:mismatch-size}; they are not identified with the
physical stopped remainder.
\end{proof}

\subsection{Counting the reasons a root can survive}

The genealogy in this subsection is the \emph{algebraic} root-family
forest. The original contact record in
Section~\ref{cell:contact-records} supplies its ordered edges. The
ownership rule below uses only that record; it does not partition a
physical force group or presume that its contact components have one
root. The coordinate Jacobians are proved later in
Lemma~\ref{hist:pinning}.

\begin{definition}[Disjoint ownership]\label{cent:ownership}
Fix the labels, their order, and an original contact record for each occupied
cell. At the top of a layer give each root label its own owner. Carry
these owners backward through the cells. In a cell form the finite
graph of physical and required virtual contacts of each original
factor's displayed generation diagram, including initial-polymer
edges. Take the disjoint union over tensor factors; contacts with
other factors remain tests and marks below, and are not used to
assign original parents. The selected
original pin edges, together with internal contact trees of its
pinned groups, connect every new label to the old array. For an
internal tree choose the lexicographically first spanning tree of
the finite contact graph. Root the resulting forest at the old
array, retaining the prescribed original parent in the singleton case
of Section~\ref{cell:contact-records}. More explicitly, process the selected edges in
their record order, with the internal-tree and label orders resolving
ties: keep an edge joining two distinct components when at least one
is rootless, and omit an edge within one component or between two
components that already contain old
labels. Orient the retained trees toward their old-label endpoints.
If needed, process the remaining required contact edges in the same
fixed order to finish this rooted forest. Connectivity guarantees
that no new label remains rootless.

Every new label inherits the owner of its old-label endpoint; an old
label never changes owner. An omitted edge between two already
owned components is still a displayed contact. If its endpoint
owners differ, mark both endpoint lives. Likewise mark both
endpoints of every other contact between different owners, whether
or not the record selected that contact as a pin. A multiple addition
marks every new label and each old label participating in its pins
or contacts. The \emph{family} of a root is the set of labels with
that root as owner. This definition changes neither a force group
nor a trajectory, and its families are disjoint by construction.
\end{definition}

If a joint root force group contains several roots, they therefore
keep different owners. A root--root encounter marks both participants;
a new connected group touching several such roots is assigned by its
chosen parent edges, while every contact to another owner marks both
sides. Thus the construction introduces no partition sum over roots
and does not silently replace a joint force group by separate flows.
Only an unmarked family will subsequently be separated, using the
absence of cross contacts and uniqueness of the prescribed flow.

For an expanded layer, also mark a displayed particle life whenever it
participates in a contact outside its scheduled isolated pair
passage, in a cross-family contact of displayed trajectories, in a
nonsingleton cut, or in a passage that crosses a required cut or fails
to finish before an interfering birth. Mark participants of a
multiple addition or endpoint-overlap insertion as well. Additional
parameter restrictions used later may mark further lives. The
\emph{participants} of a binary event are its two labels; for a
multiple event they are the added labels and every existing label
whose contact or required constraint prevents its family from being
simple. An unrelated root in the same cell is not a participant.
A particle life here is a formal label with its prescribed lifetime;
revisiting that life in another layer does not create a new life.

\begin{proposition}[Cross-layer charge count]\label{cent:charges}
Consider a term remaining after Proposition~\ref{cent:cancellation},
with no unresolved cutoff interruption. Let $H$ count the root slots
entering each expanded layer at its upper cut, including the starting
roots and residual slots, summed over those layers.
Retained inputs at time zero are terminal charges,
not additional root occurrences in $H$. Let $h_0$ count uncentered
restarting top roots, let $d$ be the number of distinct marked lives,
let $r$ be the number of residual factors, and let $i$ be the number
of retained singleton factors at time zero. With at most $L$ expanded
layers,
\begin{equation}\label{cent:charge-bound}
 H\leq Lh_0+L^2d+Lr+Li.
\end{equation}
In particular, if $H\geq\delta p$ and $h_0\leq\delta p/(2L)$, then
\begin{equation}\label{cent:extensive-charges}
 d+r+i\geq\frac{\delta}{2L^2}p.
\end{equation}
The statement is purely combinatorial and imposes no smallness on an
uncentered restarting top.
\end{proposition}
\begin{proof}
Use the disjoint ownership of Definition~\ref{cent:ownership} in each
layer. Every original singleton birth has its parent's owner. If a
family has no marked member, its births are single and every passage
is isolated, complete, and separated at its required cuts. Every
cross-owner contact would mark its endpoint in that family, so none
exists. In a joint root force group the cross potentials therefore
vanish along these trajectories, and uniqueness splits off this
family's flow. A connected detached group cannot straddle this family
and another one without a cross-owner contact. Its generation diagram
consequently splits by Proposition~\ref{cell:factorization}. This
proves that an unmarked family is precisely eligible for the simple
reference construction; it is not an assumption about the components
of the joint root group. A centered
singleton root of this family must therefore either have a residual
factor, or have a retained lower singleton descendant: otherwise its
family would cancel by Proposition~\ref{cent:cancellation}. A root
in a nonsingleton block is already marked.

For each uncharged centered root choose one retained lower descendant,
using the first one in a fixed order. Different unmarked families
have disjoint labels, so these choices are distinct. Draw an arrow
from that root occurrence to the chosen root entering the next lower
expanded layer, or end its path at the chosen initial factor if the
lower cut is time zero. The counted vertices are precisely the
entering-layer root occurrences in $H$. Root
occurrences charged to marked families have no arrow; assign each
such root one marked member of its family. This assignment is
injective within a layer because the families are disjoint.
Residual roots and uncentered restarting roots also terminate paths.
Every other vertex has one downward continuation, and every vertex
receives at most one arrow. Hence the graph is a disjoint union of
downward paths. Every path has at most $L$ counted root occurrences
and ends at one of the specified charges.

At most $h_0$ ends are restarting roots, at most $r$ are residual
roots, and at most $i$ are initial roots. At any layer the marked
family ends inject into the marked lives. One life occurs in at most
$L$ layers, so there are at most $Ld$ such ends. Multiplying the
number of ends by the maximal path length proves
\eqref{cent:charge-bound}. The last inequality follows immediately.
\end{proof}

\begin{remark}[Changing the pin forest]\label{cent:pin-forest}
Ordinary parent counting, compact-parameter estimates, and the
terminal-subtree selection use the original forest just defined.
Only the final extensive-defect estimate changes to global earliest
contact coordinates; its full $p^{Cp}$ multiplicity is then paid by
an extensive geometric gain. No gain is assigned to an arbitrary
irregular rerouting. Lemma~\ref{exp:rerouting} first controls premature
contacts outside the original parameter exceptions, and the few
remaining affected births are marked irregular.

For a singleton whose global pin is rerouted to an earlier contact
with another factor, mark its new child and its old and new parent
lives, at most three labels. That earlier contact was already a
cross-family contact in the algebraic diagram, and its later
required original-factor encounter remains a geometric defect.
For a multiple insertion retain all participant marks. This
finite-participant transfer, detailed in Proposition~\ref{exp:defects},
permits a fixed-factor adjustment of \eqref{cent:extensive-charges}.
It neither changes family ownership in the toggle nor asserts an
additional small factor for every irregular pin.
\end{remark}

For a fixed total number of labels, a surviving centered root has,
by the same path argument, an initial factor, a residual, or a marked
history. The first two contributions vanish by
\eqref{cent:initial-use} and \eqref{cent:residual-def}. Fixed-size
marked-history contributions vanish by Corollary~\ref{exp:finite},
with all independent top phase points integrated. That corollary
also supplies the finite-size majorant for the small initial and
residual factors. This uses only finite-history estimates and does
not presume that an actual-process large history is already rare.
